\documentclass[11pt]{article}
\usepackage{amsmath, amssymb, amsthm}
\usepackage{geometry}
\usepackage{hyperref}
\usepackage{booktabs}
\usepackage{listings}

\geometry{margin=1.1in}

\newtheorem{lemma}{Lemma}
\newtheorem{fact}{Fact}

\newcommand{\eps}{\varepsilon}
\newcommand{\dt}{\Delta t}
\newcommand{\dx}{\Delta x}
\newcommand{\code}[1]{\texttt{#1}}

\lstset{
  basicstyle=\ttfamily\small,
  columns=fullflexible,
  keepspaces=true,
  breaklines=true,
  frame=single,
}

\title{The ADMM Norms: $M$, $N$, and the Adjoint of $A$}
\author{Optimal Transport / Schr\"odinger Bridge Project\\
Companion to \code{nonuniform\_grid\_xy\_variables.tex} (Steps 1--5)}
\date{}

\begin{document}
\maketitle

\begin{abstract}
Continues \code{nonuniform\_grid\_xy\_variables.tex}, which pinned down (Steps 1--5,
confirmed) the two grids, $f_1,f_2$, and the linear consensus map $A(q,b)=(\bar{\mathbf
I}_tq,\bar{\mathbf I}_xb)$, and left one thing open: the inner products
$\langle\cdot,\cdot\rangle_\mathcal{X},\langle\cdot,\cdot\rangle_\mathcal{Y}$ (that
document's $\mathcal X,\mathcal Y$; this note calls their matrices $N,M$ respectively, to
match ongoing discussion). This note (a) writes down $M,N$ concretely -- the
quadrature-consistent Riemann-sum norms, already sitting in \code{pipeline\_nu.py} as a
\emph{diagnostic}, here promoted to the actual ADMM metric; (b) derives $A^*$, the
$N,M$-weighted adjoint of $A$, in closed form; (c) proves $\|A^*A\|_N\le1$
\emph{exactly}, for any grid in this family, giving an unconditional (not
grid-tuned) stability bound $\tau>\gamma$ for \S9's linearized-ADMM iteration.

A companion draft, \code{norms\_in\_ladmm.tex} (older notation: half-integer indices,
$\mathcal A,\tilde\rho$), independently derives the same weighted-average formula for the
time-averaging adjoint (its Step 5b) via direct coefficient-collection rather than the
general adjoint formula used here -- a useful cross-check, reached two different ways.
\end{abstract}

%------------------------------------------------------------------------------
\section{Setup: $M$, $N$, and $A$, restated}
\label{sec:setup}
%------------------------------------------------------------------------------

Notation and indices exactly as in \code{nonuniform\_grid\_xy\_variables.tex}: $q\in
\mathbb R^{(n_t-1)n_x}$ at $(t_j,\xi_k)$, $j=1,\ldots,n_t-1$; $b\in\mathbb
R^{n_t(n_x-1)}$ at $(\tau_i,x_k)$, $i=0,\ldots,n_t-1$; $\rho,m\in\mathbb R^{n_tn_x}$ at
$(\tau_i,\xi_k)$ (its ambient-space equation, \S1). $\dt_i:=t_{i+1}-t_i$
($i=0,\ldots,n_t-1$, the primal/cell widths); $\dt_j^{\mathrm{dual}}:=\tfrac12(\dt_{j-1}
+\dt_j)$ ($j=1,\ldots,n_t-1$, the dual-mesh cell widths -- \S8.3's $r_t$/$\bar{\mathbf
I}_t$ construction, and the ``dual grid'' discussion from the same thread).

\textbf{$M$ (collocated/$y$-grid norm):}
\begin{equation}
  \|(\rho,m)\|_M^2 = \dx\sum_{i,k}\dt_i\,\rho_{i,k}^2 \;+\; \dx\sum_{i,k}\dt_i\,m_{i,k}^2,
  \qquad
  M = \begin{pmatrix}\mathrm{diag}(\dt_i)\otimes\dx I_{n_x} & 0 \\ 0 & \mathrm{diag}(\dt_i)\otimes\dx I_{n_x}\end{pmatrix}.
\end{equation}

\textbf{$N$ (staggered/$x$-grid norm):}
\begin{equation}
  \|(q,b)\|_N^2 = \dx\sum_{j,k}\dt_j^{\mathrm{dual}}\,q_{j,k}^2 \;+\; \dx\sum_{i,k}\dt_i\,b_{i,k}^2,
  \qquad
  N = \begin{pmatrix}\mathrm{diag}(\dt_j^{\mathrm{dual}})\otimes\dx I_{n_x} & 0 \\ 0 & \mathrm{diag}(\dt_i)\otimes\dx I_{n_x-1}\end{pmatrix}.
\end{equation}

Both are literally \code{pipeline\_nu.py}'s \code{norm\_fn}/\code{dual\_norm\_fn}
(\code{\_weighted\_sq\_sum}) -- currently used only to \emph{measure} convergence, not to
\emph{run} the iteration. This note's point is to promote them to the actual metric.

\textbf{$A$ (consensus map, linear part):} $A(q,b)=(\bar{\mathbf I}_tq,\bar{\mathbf
I}_xb)$, where (\code{nonuniform\_grid\_xy\_variables.tex} §8.3, §10)
\begin{equation}
  (\bar{\mathbf I}_tq)_i = \tfrac12(q_i+q_{i+1}),\ i=0,\ldots,n_t-1\ (q_0:=q_{n_t}:=0),
  \qquad
  (\bar{\mathbf I}_xb)_k = \tfrac12(b_{k-1}+b_k),\ k=1,\ldots,n_x\ (b_0:=b_{n_x}:=0),
\end{equation}
each acting independently per space-column ($\bar{\mathbf I}_t$) or per time-row
($\bar{\mathbf I}_x$) -- $q_0,q_{n_t},b_0,b_{n_x}$ set to $0$ here since this is the
\emph{linear} part only (the affine boundary contribution is $A$'s companion, $d$, not
part of $A$ itself).

%------------------------------------------------------------------------------
\section{The weighted adjoint, in general}
\label{sec:adjoint-general}
%------------------------------------------------------------------------------

$A^*:\mathcal Y\to\mathcal X$ is defined by $\langle Ax,v\rangle_M=\langle x,A^*v
\rangle_N$ for all $x,v$. Writing $A$'s plain (array-index) matrix as $A$ too (abuse of
notation): $\langle Ax,v\rangle_M=x^TA^TMv\stackrel!=x^TNA^*v$ for all $x$, so
\begin{equation}
  \label{eq:adjoint-formula}
  A^* = N^{-1}A^TM.
\end{equation}
$N$ is SPD (diagonal, strictly positive entries), hence invertible -- always well-defined.

%------------------------------------------------------------------------------
\section{Reduction to two independent 1D problems}
\label{sec:reduction}
%------------------------------------------------------------------------------

$A$ has no cross terms between $(q\to\rho)$ and $(b\to m)$, and $M,N$ respect this same
split -- so $A^*$ splits the same way, $A^*(\rho,m)=(A_t^*\rho,\,A_x^*m)$. Within each
block, the averaging operator acts independently across many parallel ``copies'' ($n_x$
space-columns for $\bar{\mathbf I}_t$, $n_t$ time-rows for $\bar{\mathbf I}_x$), and in
both cases the weight on the axis \emph{not} being acted on is identical on the domain
and codomain side (space's $\dx$ for the $q\to\rho$ block; time's $\dt_i$, row by row, for
the $b\to m$ block). The following lemma says that's exactly the condition under which the
many copies don't interact -- the full weighted adjoint (and its norm) reduces to the
\emph{single} 1D averaging matrix's.

\begin{lemma}[Parallel copies don't change the weighted operator norm]
\label{lem:copies}
Let $P:\mathbb R^p\to\mathbb R^q$ be linear, $D_1,D_2$ SPD on $\mathbb R^p,\mathbb R^q$.
For $r$ independent copies of $P$ (i.e.\ $P\otimes I_r$ on $\mathbb R^{pr}$), weighted by
$D_1\otimes C$ on the domain and $D_2\otimes C$ on the codomain, $C:=\mathrm{diag}(c_1,
\ldots,c_r)$ \textbf{the same} per-copy scalars on both sides:
\[
  \|P\otimes I_r\|_{D_1\otimes C\to D_2\otimes C} = \|P\|_{D_1\to D_2}.
\]
\end{lemma}
\begin{proof}
The squared norm is the top generalized eigenvalue of $(D_1\otimes C)^{-1}(P\otimes
I_r)^T(D_2\otimes C)(P\otimes I_r) = (D_1^{-1}\otimes C^{-1})(P^TD_2P\otimes C) =
D_1^{-1}P^TD_2P\otimes I_r$ (the $C^{-1}C=I_r$ cancellation is exactly where ``same
scalars on both sides'' is used), whose eigenvalues are those of $D_1^{-1}P^TD_2P$, each
with multiplicity $r$.
\end{proof}

Applying Lemma~\ref{lem:copies}: the $q\to\rho$ block reduces to $\bar{\mathbf I}_t:
(\mathbb R^{n_t-1},N_t)\to(\mathbb R^{n_t},M_t)$ with $N_t=\mathrm{diag}(\dt_j^{
\mathrm{dual}})$, $M_t=\mathrm{diag}(\dt_i)$ ($c_k=\dx$, same for every space column
$k$); the $b\to m$ block reduces to $\bar{\mathbf I}_x:(\mathbb R^{n_x-1},\dx
I)\to(\mathbb R^{n_x},\dx I)$ ($c_i=\dt_i$, same for every time row $i$, since $b,m$
share the identical time weight at each row).

%------------------------------------------------------------------------------
\section{$A_t^*$, explicitly}
\label{sec:At}
%------------------------------------------------------------------------------

$A_t^*=N_t^{-1}\bar{\mathbf I}_t^TM_t$. Column $j$ of $\bar{\mathbf I}_t$ is nonzero
(value $\tfrac12$) in rows $j-1,j$ only, so
\begin{equation}
  \label{eq:At-star}
  (A_t^*v)_j \;=\; \frac{1}{\dt_j^{\mathrm{dual}}}\Bigl[\tfrac12\dt_{j-1}v_{j-1}+\tfrac12
  \dt_jv_j\Bigr] \;=\; \frac{\dt_{j-1}v_{j-1}+\dt_jv_j}{\dt_{j-1}+\dt_j}, \qquad
  j=1,\ldots,n_t-1.
\end{equation}
A \emph{width-weighted} average of $v$'s two flanking cells -- not the plain $\tfrac12,
\tfrac12$ average. It reduces to the plain average exactly when $\dt_{j-1}=\dt_j$
(locally uniform), and only then -- exactly the condition under which
\code{nonuniform\_grid\_refine\_ends.tex}'s ``\code{interp\_t\_at\_rho}: an accuracy
caveat'' section found the plain average's error vanishing. This is not a coincidence:
\eqref{eq:At-star} is \emph{forced} by $M,N$ once $A^*$ is required to be their genuine
adjoint (Fact~\ref{fact:not-coincidence} below makes this precise), the same
``structural, not approximate'' character \code{nonuniform\_grid\_refine\_ends.tex}
\S5 already claimed for the \emph{Euclidean} adjoint on the uniform grid.

\begin{fact}
\label{fact:not-coincidence}
No single scalar $c$ satisfies $A_t^*=c\cdot\bar{\mathbf I}_t^T$. $\bar{\mathbf
I}_t^T$ weighs $v_{j-1}:v_j$ in the fixed ratio $1:1$ at every $j$; $A_t^*$ weighs them
$\dt_{j-1}:\dt_j$, a ratio that changes from edge to edge on a graded grid. A fixed $1:1$
ratio scaled by any single $c$ cannot match a ratio that varies with $j$.
\end{fact}

\subsection*{Against the code}

\code{operators\_nu.py}'s \code{interp\_t\_at\_rho} -- called by \code{pipeline\_nu.py}'s
\code{At\_fn}, which builds the gradient \code{g} \emph{every} LADMM iteration, not just
inside the projection -- is
\begin{lstlisting}[language=Python]
def interp_t_at_rho(in_arr):
    return 0.5 * (in_arr[:-1] + in_arr[1:])
\end{lstlisting}
the plain average, i.e.\ the Euclidean transpose $\bar{\mathbf I}_t^T$, not
\eqref{eq:At-star}. Fact~\ref{fact:not-coincidence}: this is not off by a rescaling, it is
a genuinely different operator on any grid where consecutive $\dt_i$ differ -- which, for
Chebyshev clustering, is \emph{every} interior edge (unlike
\code{refine\_ends\_time\_grid}'s two isolated junctions).

%------------------------------------------------------------------------------
\section{$A_x^*$, explicitly}
\label{sec:Ax}
%------------------------------------------------------------------------------

Same computation with $\dx$ constant on both sides -- it cancels entirely:
\begin{equation}
  A_x^* = \bar{\mathbf I}_x^T, \qquad (A_x^*m)_k = \tfrac12(m_k+m_{k+1}).
\end{equation}
Identical to \code{operators\_nu.py}'s \code{interp\_x\_at\_m\_adj} (via
\code{interp\_x\_at\_phi} at zero BC) as coded today -- \textbf{no change needed here}.
Consistent with every earlier finding in this project: space, being uniform, never needed
correcting; only the time-direction adjoint does.

%------------------------------------------------------------------------------
\section{Bounding $\|A^*A\|_N$}
\label{sec:bound}
%------------------------------------------------------------------------------

\begin{fact}
\label{fact:AstarA}
For any bounded linear $A$ between real inner-product spaces, $\|A^*A\|_N=\|A\|_{
\mathcal X\to\mathcal Y}^2$.
\end{fact}
\begin{proof}
$\langle A^*Ax,x\rangle_N=\|Ax\|_M^2\le\|A\|^2\|x\|_N^2$; $A^*A$ is self-adjoint and PSD
w.r.t.\ $\langle\cdot,\cdot\rangle_N$, so its operator norm is its top Rayleigh quotient:
$\|A^*A\|_N=\sup_x\langle A^*Ax,x\rangle_N/\|x\|_N^2=\sup_x\|Ax\|_M^2/\|x\|_N^2=\|A\|^2$.
\end{proof}

$A$ is block-diagonal across $(q,b)\to(\rho,m)$, and $M,N$ split the same way, so
$\|A\|_{\mathcal X\to\mathcal Y}=\max\bigl(\|\bar{\mathbf I}_t\|_{N_t\to M_t},\,\|
\bar{\mathbf I}_x\|_{\dx I\to\dx I}\bigr)$ (a Pythagorean argument: $\|Ax\|_M^2=\|
\bar{\mathbf I}_tq\|_{M_t}^2+\|\bar{\mathbf I}_xb\|_{M_x}^2\le\max(\cdot)^2\bigl(\|q\|_{
N_t}^2+\|b\|_{N_x}^2\bigr)=\max(\cdot)^2\|x\|_N^2$, tight by concentrating $x$ on
whichever block attains the max).

\begin{lemma}[Both blocks are $1$-Lipschitz]
\label{lem:1lip}
$\|\bar{\mathbf I}_t\|_{N_t\to M_t}\le1$ and $\|\bar{\mathbf I}_x\|_{\dx I\to\dx I}\le1$,
for \emph{any} $\{\dt_i\}$ (uniform or not) and any $\dx$.
\end{lemma}
\begin{proof}
Write $u_j:=q_j$ ($j=1,\ldots,n_t-1$, $u_0:=u_{n_t}:=0$ -- the linear-part convention of
\S\ref{sec:setup}), $v_i:=(\bar{\mathbf I}_tq)_i=\tfrac12(u_i+u_{i+1})$. By
Cauchy--Schwarz, $v_i^2\le\tfrac12(u_i^2+u_{i+1}^2)$, so
\[
  \sum_{i=0}^{n_t-1}\dt_iv_i^2 \;\le\; \tfrac12\sum_i\dt_i(u_i^2+u_{i+1}^2)
  \;=\; \tfrac12\sum_{j=1}^{n_t-1}(\dt_{j-1}+\dt_j)u_j^2
  \;=\; \sum_{j=1}^{n_t-1}\dt_j^{\mathrm{dual}}u_j^2,
\]
i.e.\ $\|\bar{\mathbf I}_tq\|_{M_t}^2\le\|q\|_{N_t}^2$ -- using only
$\dt_j^{\mathrm{dual}}:=\tfrac12(\dt_{j-1}+\dt_j)$'s definition, nothing about
uniformity. The $\bar{\mathbf I}_x$ case is the same argument with every weight equal to
$\dx$ (the special case $\dt_{j-1}=\dt_j=\dx$ throughout).
\end{proof}
Essentially tight, not just a safe bound: a constant $q\equiv c$ pushes the ratio toward
$1$ as boundary effects wash out.

\begin{equation}
  \label{eq:final-bound}
  \boxed{\;\|A^*A\|_N \;\le\; 1,\quad\text{for any }n_t,n_x,\text{ any clustering.}\;}
\end{equation}

%------------------------------------------------------------------------------
\section{Consequence: an unconditional stability bound}
\label{sec:consequence}
%------------------------------------------------------------------------------

\code{nonuniform\_grid\_xy\_variables.tex} \S9's linearized-ADMM iteration needs
$\tau>\gamma\|A\|_{\mathcal X\to\mathcal Y}^2$ for convergence (Fang--He--Liu--Yuan
2015). By \eqref{eq:final-bound}, $\|A\|^2\le1$ \emph{exactly}, on every grid in this
family -- so \textbf{$\tau>\gamma$ suffices unconditionally}: no per-grid retuning, no
asymptotic argument, a genuine equality-free bound holding at every finite $n_t,n_x$ and
every clustering strength. This is the first piece of the ``grid-independent ADMM''
goal delivered as a proof rather than an aspiration -- contingent, of course, on actually
switching the iteration's $A^*$ (\S\ref{sec:At}--\S\ref{sec:Ax}) and the projection's
adjoint machinery over to $M,N$, which is not yet done.

%------------------------------------------------------------------------------
\section{Empirical note: why $R^*$ can't wait}
\label{sec:empirical}
%------------------------------------------------------------------------------

Wiring $A_t^*$ into \code{pipeline\_nu.py}'s \code{At\_fn} alone -- leaving
\code{projection\_nu.py}'s internal adjoint on the Euclidean metric -- was tried and
reverted: it satisfies the adjoint identity (\S\ref{sec:At}, verified) and reduces to the
old code on a uniform grid, but on an actual clustered grid
(\code{sb\_gaussian\_nonuniform.py}, $n_t=n_x=64$) it made the LADMM iteration diverge
(residuals $\sim10^{20}$), even though the projection still exactly enforced $R=0$
throughout. Simplifying the $y$-update alone (\code{xy\_variables.tex} \S11.5, eq.\ 24)
is stable but converges to a visibly worse answer than the unmodified solver
($\sim$25$\times$ worse error against a Sinkhorn reference). The gradient step's adjoint
and the projection's internal adjoint have to agree on one consistent metric for the
linearized-ADMM majorization argument to hold -- these pieces are not independently
swappable. Both reverts are documented in \code{prox\_nu.py}/\code{pipeline\_nu.py}'s
current docstrings.

%------------------------------------------------------------------------------
\section{$R$'s weighted adjoint (the FP projection's remaining piece)}
\label{sec:Rstar}
%------------------------------------------------------------------------------

\textbf{This section is a derivation to check, not a finished result.} \S\ref{sec:empirical}
is why this piece must be right, and consistent with $A_t^*$, before anything here is
implemented again. Unlike \S\ref{sec:At}'s $A_t^*$, nothing below has been checked
against a numerical adjoint identity yet.

\subsection{$R$'s Euclidean transpose, per building block}

$R(Q,B) = \mathbf D_tQ + B\mathbf D_x^T - \varepsilon(\bar{\mathbf I}_tQ)\mathbf L_x^T$
(\code{nonuniform\_grid\_xy\_variables.tex} \S8.3--8.4's matrices, $Q,B$ the matrix forms
of $q,b$: $\mathbf D_t,\bar{\mathbf I}_t$ act by left-multiplication -- per space column
-- and $\mathbf D_x,\mathbf L_x$ by right-multiplication -- per time row, exactly
\S\ref{sec:reduction}'s convention). Taking the plain Euclidean (Frobenius) adjoint of
each term ($\mathbf L_x=\mathbf L_x^T$, symmetric):
\begin{equation}
  \label{eq:R-transpose}
  (R^T\Phi)_q = \mathbf D_t^T\Phi - \varepsilon\,\bar{\mathbf I}_t^T(\Phi\mathbf L_x),
  \qquad
  (R^T\Phi)_b = \Phi\mathbf D_x.
\end{equation}
This matches \code{projection\_nu.py} as coded today (up to the module's own
already-verified sign convention, \code{deriv\_t\_at\_rho}$=-\mathbf D_t^T$ --
see \code{validate\_nu.py}'s \code{check\_deriv\_t\_adjoint}):
\begin{lstlisting}[language=Python]
rho_correction = deriv_t_at_rho(phi) + vareps * interp_t_at_rho(deriv_x_at_phi(deriv_x_at_m(phi)))  # (R^T phi)_q
mx_correction  = deriv_x_at_m(phi)                                                                     # (R^T phi)_b
\end{lstlisting}

\subsection{$R^*=N^{-1}R^TW$, for \emph{any} invertible $W$ -- the choice is irrelevant}

\begin{lemma}[The codomain weight cancels]
\label{lem:W-cancels}
For any invertible $W$ (not just $W=I$), $R^*:=N^{-1}R^TW$ gives the \emph{same}
projection formula as $W=I$:
\[
  v - R^*(RR^*)^{-1}(Rv-d) \;=\; v - N^{-1}R^T(RN^{-1}R^T)^{-1}(Rv-d), \qquad \forall\,W.
\]
\end{lemma}
\begin{proof}
Write $K:=RN^{-1}R^T$. Then $RR^*=RN^{-1}R^TW=KW$, so $(RR^*)^{-1}=(KW)^{-1}=W^{-1}K^{-1}$,
and $R^*(RR^*)^{-1}=N^{-1}R^TW\cdot W^{-1}K^{-1}=N^{-1}R^TK^{-1}$ -- the $W,W^{-1}$ cancel
identically, for any invertible $W$.
\end{proof}
So $R$'s codomain (the residual/dual-potential space, where $\phi$ lives) does not need
its own weight assigned at all -- not because it \emph{has} none (it could just as well be
$M$, since $\phi$ lives at the same points $\rho$ does), but because
Lemma~\ref{lem:W-cancels} makes the choice provably immaterial to the projection this
computes. Skipping it (equivalently: computing with $W=I$, the cheapest representative of
every equally-valid choice) is not an approximation or a simplifying assumption:
\begin{equation}
  \label{eq:R-star}
  R^* = N^{-1}R^T, \qquad
  (R^*\Phi)_q = N_q^{-1}\Bigl[\mathbf D_t^T\Phi - \varepsilon\,\bar{\mathbf I}_t^T(\Phi
  \mathbf L_x)\Bigr], \qquad
  (R^*\Phi)_b = N_b^{-1}\bigl[\Phi\mathbf D_x\bigr],
\end{equation}
$N_q^{-1}$ applied once to the whole $q$-correction (not to each piece before combining
-- linear, so the two give the same answer either way).

\subsection{The $N^{-1/2}$ connection}

Checking this against the whitening view directly answers ``isn't it just $N^{-1/2}$'':
$\tilde R:=RN^{-1/2}$, $u:=N^{1/2}v$. Solve $\tilde R\tilde R^T\phi=\tilde Ru\,(=Rv$, since
$\tilde Ru=RN^{-1/2}N^{1/2}v=Rv)$, then $x^*=N^{-1/2}(u-\tilde R^T\phi)=v-N^{-1}R^T\phi$.
So: \emph{one} factor of $N^{-1/2}$ modifies $R$ itself (giving $\tilde R$), and it
appears \emph{twice} once $\tilde R\tilde R^T$ is formed -- contributing $N^{-1}$
overall, matching \eqref{eq:R-star} exactly. Not two independent $N^{-1/2}$'s; one,
appearing on both sides of the assembled system.

\subsection{What \code{\_probe\_composite\_ops} computes today}

The module doesn't use \eqref{eq:R-transpose} in closed form -- it probes with unit
vectors $v=e_j$ (module docstring: ``rather than hand-rederive the closed-form
coefficients for a variable \code{dt\_vec}... this module builds the per-mode operator
by literal substitution''):
\begin{lstlisting}[language=Python]
drho = ops.deriv_t_at_rho(v)      # (D_t^T v)-piece, sign per the module's convention
irho = ops.interp_t_at_rho(v)      # (I_t^T v)-piece
M0[:,j] = -ops.deriv_t_at_phi(drho, 0, 0)
M1[:,j] = vareps*(ops.deriv_t_at_phi(irho,0,0) - ops.interp_t_at_phi(drho,0,0)) + v[:,0]
M2[:,j] = ops.interp_t_at_phi(irho, 0, 0)
\end{lstlisting}
Cross-checked against \code{projection.py}'s closed form (uniform grid, $M0,M1,M2$'s
role there): $M0$ is the ($q$-side, time-only) $\mathbf D_t\mathbf D_t^T$ contribution
($2/\dt^2$ on the diagonal, matching \code{projection.py}'s \code{D[1:-1]}); the bare
``\code{+ v[:,0]}'' in $M1$ is the ($b$-side) $\mathbf D_x\mathbf D_x^T$ contribution --
DCT-diagonalized to a pure scalar eigenvalue $l_{x,k}$ (\code{projection.py}: plain
``\code{+ lx}'' in \code{D[1:-1]}, no other factor), which is why it shows up as bare
\code{v} (identity in time) rather than any $D_t$-built stencil -- there is no $\varepsilon$
or $D_t$ dependence in that specific piece to probe; the rest of $M1$ and all of $M2$ are
the diffusion term's self- and cross-interactions.

\subsection{Solving the projection: the space eigenvalues survive untouched}
\label{sec:space-survives}

\eqref{eq:R-star} gives $R^*$; the projection itself needs $R(R^*\Phi)$ ($d=0$ for $R$
specifically, the FP-feasible set being $\{Rx=0\}$, \code{xy\_variables.tex} \S8.4).
Substitute \eqref{eq:R-star}'s
$Q^*:=(R^*\Phi)_q,\,B^*:=(R^*\Phi)_b$ into $R(Q,B)=\mathbf D_tQ+B\mathbf D_x^T-
\varepsilon(\bar{\mathbf I}_tQ)\mathbf L_x^T$ and expand every term:
\begin{align}
  \mathbf D_tQ^* &= \underbrace{\bigl[\mathbf D_t\mathbf N_q^{-1}\mathbf D_t^T\bigr]}_{\text{time-only}}\Phi
    \;-\;\varepsilon\underbrace{\bigl[\mathbf D_t\mathbf N_q^{-1}\bar{\mathbf I}_t^T\bigr]}_{\text{time-only}}(\Phi\mathbf L_x), \label{eq:term1}\\
  B^*\mathbf D_x^T &= \underbrace{\mathbf N_b^{-1}}_{\text{time-only}}\,\Phi(\mathbf D_x\mathbf D_x^T), \label{eq:term2}\\
  -\varepsilon(\bar{\mathbf I}_tQ^*)\mathbf L_x^T &=
    -\varepsilon\underbrace{\bigl[\bar{\mathbf I}_t\mathbf N_q^{-1}\mathbf D_t^T\bigr]}_{\text{time-only}}(\Phi\mathbf L_x)
    \;+\;\varepsilon^2\underbrace{\bigl[\bar{\mathbf I}_t\mathbf N_q^{-1}\bar{\mathbf I}_t^T\bigr]}_{\text{time-only}}(\Phi\mathbf L_x^2). \label{eq:term3}
\end{align}
(Using $\mathbf N_q^{-1},\mathbf N_b^{-1}$ for the diagonal matrices $N_q,N_b$ act as --
\S\ref{sec:reduction}'s left/right-multiplication convention; $\mathbf L_x=\mathbf L_x^T$.)
\textbf{Every single term factors as (a time-only matrix, the same at every spatial
mode) times $\Phi$, $\Phi\mathbf L_x$, or $\Phi\mathbf L_x^2$} -- and
$\mathbf D_x\mathbf D_x^T,\mathbf L_x,\mathbf L_x^2$ are exactly the operators the
\emph{current, unweighted} code already DCT-diagonalizes, with exactly the same
eigenvalues $\lambda_{x,k},\lambda_{x,k}^2$ it already uses
(\code{problem.lambda\_x}) -- $\mathbf N_q^{-1},\mathbf N_b^{-1}$ never appear inside a
space-direction factor, only outside/beside one, because \eqref{eq:term1}--\eqref{eq:term3}
never multiply a time matrix and a space matrix from the same side of $\Phi$. \textbf{This
is the ``space part stays exactly the same'' you asked me to check: it does, term by
term, not just plausibly.} After DCT-in-$x$, mode $k$'s reduced system is still exactly
\begin{equation}
  \label{eq:Tk-weighted}
  \widetilde T_k = \widetilde M0 + \lambda_{x,k}\widetilde M1 + \lambda_{x,k}^2\widetilde M2,
\end{equation}
the same quadratic-in-$\lambda_{x,k}$ shape \code{precomp\_banded\_proj\_nu} already
assembles -- only the (still exactly tridiagonal, since every piece above is built from
compositions of the bidiagonal $\mathbf D_t,\bar{\mathbf I}_t$) time-direction matrices
change:
\begin{equation}
  \label{eq:M-weighted}
  \widetilde M0 = \mathbf D_t\mathbf N_q^{-1}\mathbf D_t^T, \qquad
  \widetilde M2 = \bar{\mathbf I}_t\mathbf N_q^{-1}\bar{\mathbf I}_t^T, \qquad
  \widetilde M1 = \pm\Bigl[\mathbf D_t\mathbf N_q^{-1}\bar{\mathbf I}_t^T+\bar{\mathbf I}_t
  \mathbf N_q^{-1}\mathbf D_t^T\Bigr] + \mathbf N_b^{-1},
\end{equation}
$\widetilde M0,\widetilde M2$'s signs match \eqref{eq:term1}/\eqref{eq:term3} directly
(both unambiguous); $\widetilde M1$'s $\pm$ is deliberately left open here -- see below.

\subsection*{Where I got stuck, honestly}

I tried to pin down $\widetilde M1$'s sign by hand, against the \emph{unweighted}
($\mathbf N_q^{-1}\to I$) case, two ways, and they disagreed: my clean derivation
above gives $-\varepsilon[\mathbf D_t\bar{\mathbf I}_t^T+\bar{\mathbf I}_t\mathbf D_t^T]$;
reading the sign directly off the current code's \code{M1 = vareps*(deriv\_t\_at\_phi(irho)
- interp\_t\_at\_phi(drho)) + v} (substituting \code{drho}$=-\mathbf D_t^Tv$,
\code{irho}$=\bar{\mathbf I}_t^Tv$, the module's established sign convention) gives
$+\varepsilon[\ldots]$. Chasing the discrepancy further, I found the module's own
docstring already flags exactly this kind of trap: \code{deriv\_x\_at\_phi(deriv\_x\_at\_m(.))}
(part of what feeds $\widetilde M1$'s $\lambda_x$-linear structure) has DCT eigenvalue
\emph{$-\lambda_{x,k}$}, not $+\lambda_{x,k}$ -- ``verified numerically,'' per the module
docstring, precisely \emph{because} hand-tracking it is error-prone. This is exactly the
situation \code{\_probe\_composite\_ops} was built to avoid by construction (module
docstring: ``easy to get a sign or index wrong with no way to self-check the
algebra''), and I don't want to guess past it.

\textbf{What doesn't depend on resolving this:} \emph{which} building blocks
$N_q^{-1},N_b^{-1}$ multiply -- \code{drho}, \code{irho}, and the bare \code{+v} term,
respectively, exactly \eqref{eq:M-weighted}'s structure -- is settled by
\eqref{eq:term1}--\eqref{eq:term3} regardless of sign, since rescaling an already-correct
quantity by a positive diagonal matrix cannot flip which pieces combine. So the code
change is still exactly this (unchanged from before):
\begin{lstlisting}[language=Python]
drho = ops.deriv_t_at_rho(v) / dt_dual[:, None]     # N_q^{-1}, once
irho = ops.interp_t_at_rho(v) / dt_dual[:, None]     # N_q^{-1}, once
M0[:,j] = -ops.deriv_t_at_phi(drho, 0, 0)
M1[:,j] = vareps*(ops.deriv_t_at_phi(irho,0,0) - ops.interp_t_at_phi(drho,0,0)) \
          + v[:,0] / dt_vec                          # N_b^{-1}, NOT N_q^{-1}
M2[:,j] = ops.interp_t_at_phi(irho, 0, 0)
\end{lstlisting}
plus, separately (\eqref{eq:R-star}'s $b$-block -- applied directly where \code{mx\_out}
is built, outside \code{\_probe\_composite\_ops} entirely):
\code{mx\_out = psi + dphi\_dx / dt\_vec[:,None]}. What this section adds is
not a sign, but the reason not to hand-derive one: exactly like the rest of this module,
the actual $\pm$ in \eqref{eq:M-weighted} should come from running this code and checking
it against $R,R^*$'s adjoint identity numerically (\S\ref{sec:Rstar}'s closing
``To verify''), not from further hand algebra.

\subsection*{To verify -- before touching \code{projection\_nu.py} again}

\begin{enumerate}
  \item Does \eqref{eq:R-transpose}--\eqref{eq:R-star} look right?
  \item The proposed fix above -- in particular, $N_q^{-1}$ on \code{drho}/\code{irho}
        but a \emph{different}, $N_b^{-1}$ on $M1$'s bare \code{+v} term -- is that the
        right split, or is there a piece still missing?
  \item Given last time: \textbf{before any code changes}, write a numerical
        adjoint-identity check for $R^*$ against $R$ (mirroring
        \code{check\_At\_star\_is\_NM\_weighted\_adjoint}, \S\ref{sec:At}'s pattern) --
        agreed this is the right gate before implementation, given the empirical
        divergence \S\ref{sec:empirical} found last time from skipping exactly that step
        for the projection?
\end{enumerate}
\end{document}
